\chapter{流形上的积分学}\label{appx:form}

\section{外微分}
有了流形上的微分学，自然要有相应的积分学。
观察 $\R^n$ 上的多元 Riemann 积分，如 $\int P\,\d x+Q\,\d y+R\,\d z,\int P\,\d x\,\d y+Q\,\d y\,\d z+R\,\d z\,\d x$，
被积内容总按微元的组合来求和，故可视作微分形式，系数则是其坐标分量。如此，积分便可写成 $\int \omega$，微元间的乘法可理解为楔积。
再研究形式的微分。函数 $f$ 的微分 $\d f$ 给出 1-形式，该操作增添了一个微元。推广到形式的微分无非是作用于形式的系数并添上楔积。

\begin{definition}
    流形上的形式场称为\textbf{微分形式场}。常取对偶坐标基将 $l$-微分形式 $\omega$ 表为
    \eq{
    \omega=\sum_{\mu_1<\cdots<\mu_l}\omega_{\mu_1\cdots\mu_l} \d x^{\mu_1}\wedge\cdots\wedge \d x^{\mu_l}=\frac{1}{l!}\omega_{\mu_1\cdots\mu_l} \d x^{\mu_1}\wedge\cdots\wedge \d x^{\mu_l}.
    }
\end{definition}

\begin{definition}
    \textbf{外微分}(exterior derivative)是映射 $\d:\Lambda_l(V)\to\Lambda_{l+1}(V)$ 使
    \eq{
        \d\omega :=\sum_{\mu_1<\cdots<\mu_l}\d\omega_{\mu_1\cdots\mu_l}\wedge\d x^{\mu_{1}}\wedge\cdots\wedge\d x^{\mu_{l}}=\frac{1}{l!}\d\omega_{\nu_1\cdots\nu_l}\wedge\d x^{\nu_{1}}\wedge\cdots\wedge\d x^{\nu_{l}}.
    }
\end{definition}

\begin{theorem}
    任取坐标系有 $(\d\omega)_{\sigma\mu_1\cdots\mu_l}=(l+1)\omega_{[\mu_1\cdots\mu_l,\sigma]}$。
\end{theorem}

\begin{proof}由楔积定义知 $e^{\nu_1}\wedge\cdots\wedge e^{\nu_n}(e_{\mu_1},\cdots,e_{\mu_n})=n!\delta_{\mu_1}^{[\nu_1}\cdots\delta_{\mu_n}^{\nu_n]}$，则
    \begin{align*}
        \d\omega(\p_{\sigma},\p_{\mu_1},\cdots,\p_{\mu_l})&=\frac{1}{l!}\omega_{\nu_1\cdots\nu_l,\lambda}\d x^\lambda\wedge\d x^{\nu_{1}}\wedge\cdots\wedge\d x^{\nu_{l}}(\p_{\sigma},\p_{\mu_1},\cdots,\p_{\mu_l})\\
        &=\frac{(l+1)!}{l!} \omega_{[\nu_1\cdots\nu_l,\lambda]} \delta_{\sigma}^{\lambda}\delta_{\mu_1}^{\nu_1}\cdots\delta_{\mu_l}^{\nu_l}=(l+1)\omega_{[\mu_1\cdots\mu_l,\sigma]}.\qedhere
    \end{align*}
\end{proof}

\begin{theorem}[Leibniz 律]$\omega$ 为 $k$-形式：
    $\d(\omega\wedge\mu)=\d\omega\wedge\mu+(-1)^k\omega\wedge\d\mu$。
\end{theorem}

\begin{theorem}[Poincaré 引理] $\d^2=0$.
\end{theorem}
\begin{proof} $(\d\d\omega)_{\nu\sigma\mu_1\cdots\mu_l}=(l+2)(l+1)\omega_{[[\mu_1\cdots\mu_l,\sigma],\nu]}=(l+2)(l+1)\omega_{[\mu_1\cdots\mu_l,(\sigma\nu)]}=0$。
\end{proof}

\begin{definition}
    $l$-形式 $\omega$ 称\textbf{闭的}(closed) $\Leftrightarrow\d\omega=0$；$\omega$ 称\textbf{恰当的}(exact) $\Leftrightarrow \exists (l-1)$-形式 $\mu$ 使 $\omega=\d\mu$。
\end{definition}
依 Poincaré 引理，恰当形式必为闭形式。反之欲成立，须要求坐标域$X\subset M$ 第二种 de Rham 上同调 $H^{2}(X;\R)$ 平凡，见 \S\ref{appx:Alge-Topo}。形象理解为区域是任意的单连通流形，例如 $\R^{n}$、开球及更一般的“星形式”开区域等。对一般流形而言，往往局域 $X$ 内亦成立，即逆命题至少在处处点的邻域上成立。

\begin{definition}
    $\delta={\star\d\star}$ 称为\textbf{余微分}(codifferential)，有时也记作 $\d^*$。
\end{definition}
\begin{theorem}
    易证 $\delta^2=0$。
\end{theorem}

\begin{example}[矢量分析]
    梯度、旋度、散度、Laplace 算子\footnote{$\d \delta+\delta \d$ 强调为 \textbf{Laplace-de Rham 算子}。闵氏时空 d'Alembert 算符强调为 \textbf{Hodge-Laplace 算子}。}的外微分语言为
    \[\grad f=(\d f)^\sharp,\quad
    \curl\bm A=({\star\d} A_\flat)^\sharp,\quad
    \div\bm A=\delta A_\flat,\quad
    \nabla^2=\d \delta+\delta \d=(\d+\delta)^2.\]
    由此，易从 $\d^2=0$ 证明 $\curl\grad f=\bm 0$和$\div\curl\bm A=0$。
    
    对高维流形，旋度可直接用形式表述而不必 Hodge 对偶，比如四维矢量场 $V^\mu$ 的旋度就是 $2\nabla_{[\mu} V_{\nu]}$，无挠性使它在任意坐标系等价于 $2\p_{[\mu} V_{\nu]}$；二阶反称张量 $F_{\mu\nu}$ 的旋度则记作 $3\nabla_{[\mu}F_{\nu\lambda]}$。
\end{example}

\begin{example}
    电磁张量 $F_{\mu\nu}$ 正是 $\R^{3+1}$ 上的 $2$-形式。\textbf{电流密度形式} $J$ 为
\eq{J=\rho U_\flat=-J^0 \,\mathrm{d} x^{0}+J^1 \d{x^1}+J^2 \d{x^2}+J^3 \d{x^3},}
$\rho$ 为固有电荷密度。进而 Maxwell 方程组的外微分语言为
\eq{\d F=0,\quad \delta F=4\pi J.}
可见 $F$ 是闭形式。
第二式又可表为 $\mathrm{d}F^*=4\pi{\star J}$，其中 $F^*:={\star}F$ 称为\textbf{对偶电磁张量}。
某系所测电磁场为 $E=F(\oo,\p_0),B=-F^*(\oo,\p_0)$，因而 $F$ 在该系可表为
\begin{align*}F&= E^{1} \mathrm{d} x^{1} \wedge \mathrm{d} x^{0}+E^{2} \mathrm{d} x^{2} \wedge \mathrm{d} x^{0}+E^{3} \mathrm{d} x^{3} \wedge \mathrm{d} x^{0} \\&+ B^{3} \mathrm{d} x^{1} \wedge \mathrm{d} x^{2}+B^{2} \mathrm{d} x^{3} \wedge \mathrm{d} x^{1}+B^{1} \mathrm{d} x^{2} \wedge \mathrm{d} x^{3}.\end{align*}
由 $\d F=0$，一定存在 $1$-形式 $A$ 满足
\eq{F=\d A.}
$A$ 正是 $4$-势。$A$ 当然不唯一，因为 $\d^{2}=0$，$A+\mathrm{d} f$ 亦满足。但可证明存在 $A$ 满足 Lorenz 规范 $\delta A=0$，从而只需求解 $4\pi J=\delta\d A=\square A$，其中 $\square=(\d+\delta)^2$。
下面证之。从任一满足 $\d A=F$ 的势出发，令 $A'=A+\d f$。由 $\d^2=0$，仍有 $\d A'=F$。又因标量场满足 $\delta f=0$，故 $\delta A'=\delta A+\delta\d f=\delta A+\square f$。
在闵氏时空上，非齐次波动方程 $\square f=-\delta A$ 的光滑初值问题可解，故可取 $f$ 使 $\delta A'=0$。于是 $\square A'=(\d\delta+\delta\d)A'=\delta F=4\pi J$。
\end{example}

\begin{theorem}[Cartan]
    设 $X$ 为光滑切矢场，则在微分形式上有
    \eq{
        \Lie_X=\d i_X+i_X\d,\qquad [\Lie_X,\d]=0.
    }
    前者称为\textbf{Cartan 同伦式}或\textbf{魔术公式}。
\end{theorem}
\begin{proof}
    记 $D:=\d i_X+i_X\d$。对 $k$-形式 $\alpha$，内乘满足
    \eq{i_X(\alpha\wedge\beta)=(i_X\alpha)\wedge\beta+(-1)^k\alpha\wedge i_X\beta.}
    结合外微分的 Leibniz 律，展开后交叉项相消，得
    \eq{D(\alpha\wedge\beta)=(D\alpha)\wedge\beta+\alpha\wedge D\beta.}
    故 $D$ 与 $\Lie_X$ 都满足通常的乘积法则。又由前文，
    \eq{Df=i_X\d f=Xf=\Lie_Xf,\qquad D(\d x^i)=\d X^i=\Lie_X(\d x^i).}
    任意形式局部都是 $f\,\d x^{i_1}\wedge\cdots\wedge\d x^{i_l}$ 之和，故由线性性及乘积法则得 $D=\Lie_X$。最后由 $\d^2=0$，
    \eq{\Lie_X\d=\d i_X\d=\d\Lie_X.\qedhere}
\end{proof}

\section{定向与带边流形}\label{sec:orientation-boundary}

\S\ref{sec:tensor-density} 已定义线性空间的体元与定向。对流形，每个切空间都可单独定向，问题在于能否随点光滑地选取。

\begin{definition}
    $n$ 维流形 $M$ 上处处非零的光滑 $n$-形式 $\epsilon$ 称为\textbf{体元场}，简称\textbf{体元}。存在体元的流形称为\textbf{可定向的}(orientable)。沿用线性空间的约定，$\epsilon_1\sim\epsilon_2$ 指 $\epsilon_1=h\epsilon_2$，其中 $h$ 为处处正的光滑函数；其等价类称为 $M$ 的\textbf{定向}。选定定向的流形称为\textbf{定向的}(oriented)。
\end{definition}

正定体元、右手基等术语均逐点沿用。坐标基为右手基的坐标系称为\textbf{右手系}；在其中，正定体元表为
\eq{\epsilon=f\,\d x^1\wedge\cdots\wedge\d x^n,\qquad f>0.}
若 $M$ 连通且可定向，任意两体元之比为连续非零函数，故符号恒定，于是恰有两个定向；不连通时，各连通分支可独立选择定向。

\begin{theorem}
    $n\geqslant1$ 维无边流形 $M$ 可定向当且仅当存在图册，使任意坐标变换的 Jacobi 行列式处处为正。
\end{theorem}
\begin{proof}
    若有体元，缩小坐标域并在必要时反转一个坐标，即可使其成为右手系。由
    \eq{\d x'^1\wedge\cdots\wedge\d x'^n
        =\det\!\left(\pdv{x'^i}{x^j}\right)\d x^1\wedge\cdots\wedge\d x^n,}
    两右手系间的 Jacobi 行列式为正。

    反之，取从属于该图册的光滑单位分解 $\{\rho_\alpha\}$，将各 $\rho_\alpha\,\d x_\alpha^1\wedge\cdots\wedge\d x_\alpha^n$ 在坐标域外延拓为零。它们局部有限的和光滑，且在每点都是同向体元的非负组合，至少一项非零，故为体元。
\end{proof}

\begin{example}
    $\R^n$ 的标准定向由 $\d x^1\wedge\cdots\wedge\d x^n$ 给出。M\"obius 带则不可定向：将矩形按 $(0,t)\sim(1,-t)$ 粘合，若存在体元，其拉回为 $f(s,t)\,\d s\wedge\d t$，其中 $f$ 连续且处处非零。粘合要求 $f(1,0)=-f(0,0)$，与 $f(s,0)$ 符号恒定矛盾。
\end{example}

\begin{figure}[h!]
    \centering
\tikzset{every picture/.style={line width=0.75pt}} %set default line width to 0.75pt        

\begin{tikzpicture}[x=0.75pt,y=0.75pt,yscale=1,xscale=1]
%uncomment if require: \path (0,300); %set diagram left start at 0, and has height of 300

%Curve Lines [id:da7826517172347827] 
\draw    (3.65,18.88) .. controls (17.2,36.58) and (48,50.16) .. (72.8,51.99) ;
\draw [shift={(39.71,43.88)}, rotate = 205.45] [fill={rgb, 255:red, 0; green, 0; blue, 0 }  ][line width=0.08]  [draw opacity=0] (8.4,-2.1) -- (0,0) -- (8.4,2.1) -- cycle    ;
%Curve Lines [id:da8566552148207576] 
\draw    (221.73,12.59) .. controls (241.17,-1.67) and (128.44,8) .. (140.66,18.19) ;
%Curve Lines [id:da14270801999584393] 
\draw  [dash pattern={on 3pt off 3pt on 3pt off 3pt}]  (183.2,21.91) .. controls (174,25.58) and (149.2,26.68) .. (140.66,18.19) ;
%Curve Lines [id:da5926224666835718] 
\draw  [dash pattern={on 3pt off 3pt on 3pt off 3pt}]  (221.73,12.59) .. controls (206.74,22.26) and (214.4,6.87) .. (183.2,21.91) ;
%Curve Lines [id:da4918795365326445] 
\draw    (223.95,9.53) .. controls (217.84,52.82) and (166.2,65.55) .. (140.66,18.19) ;
%Straight Lines [id:da406966878688231] 
\draw    (163.2,42.45) -- (149.8,65.67) ;
\draw [shift={(148.8,67.4)}, rotate = 299.99] [fill={rgb, 255:red, 0; green, 0; blue, 0 }  ][line width=0.08]  [draw opacity=0] (7.2,-1.8) -- (0,0) -- (7.2,1.8) -- cycle    ;

\end{tikzpicture}
    \caption{曲线、曲面的定向}
\end{figure}


为讨论积分区域的边界，将坐标模型从欧氏空间推广到闭半空间
\eq{\mathbb H^n:=\{(x^1,\cdots,x^n)\in\R^n:x^n\geqslant0\}.}
半空间上的光滑映射，指在每点附近可延拓为欧氏开集上的光滑映射。

\begin{definition}
    仿紧 $T_2$ 空间 $M$ 若有图册，使各坐标域同胚于 $\mathbb H^n$ 的相对开集，且坐标变换及其逆均光滑，则赋以相应最大图册后称为 $n$ 维\textbf{光滑带边流形}(smooth manifold with boundary)。坐标满足 $x^n=0$ 的点称为\textbf{边界点}，全体记为 $\partial M$；其余点组成\textbf{内部} $M^\circ:=M\backslash\partial M$。
\end{definition}

\begin{remark}
    内点、边界点的区分与坐标无关。否则某坐标变换将内点送到边界点，其末坐标在该内点处取局部最小值，导数为零，与坐标变换的微分可逆矛盾。带边流形允许 $\partial M=\mt$，此时就是前文的无边流形。上述半空间取 $n\geqslant1$；零维流形规定无边。
\end{remark}

\begin{theorem}\label{thm:boundary-no-boundary}
    $n$ 维带边流形的非空边界是 $n-1$ 维无边流形，故 $\partial(\partial M)=\mt$。
\end{theorem}
\begin{proof}
    边界是 $M$ 的闭子集，故仍为仿紧 $T_2$ 空间。将边界坐标系限制到 $x^n=0$，再去掉末坐标，便得到取值于 $\R^{n-1}$ 开集的坐标系；其坐标变换由原坐标变换限制而来，仍双向光滑。
\end{proof}

边界点处仍有 $\dim T_pM=n$；光滑函数的局部延拓给出通常的坐标导数。通过包含映射的推前，将 $T_p(\partial M)$ 看作 $T_pM$ 的子空间，则
\eq{T_p(\partial M)=\ker(\d x^n)_p
    =\operatorname{span}\{\p_1|_p,\cdots,\p_{n-1}|_p\}.}
体元与定向的定义也照常适用。光滑单位分解的构造只需将欧氏空间的隆起函数限制到半空间，故仍成立。

\begin{remark}
    此处的边界是流形自身的概念，应与子集在环境空间中的拓扑边界区分。例如，球面作为流形无边，但作为欧氏空间的子集，其拓扑边界是自身。\textbf{闭流形}指紧且无边的流形，不能将“紧”换成“有界”。这里讨论光滑边界，不包含指定嵌入下的棱角。
\end{remark}

\begin{example}
    若无边流形 $M$ 的闭子集 $D$ 在每个拓扑边界点附近都有坐标系 $(O,\psi)$，使
    \eq{\psi[D\cap O]=\psi[O]\cap\mathbb H^n,}
    则 $D$ 在诱导拓扑及限制图册下是与 $M$ 同维的带边流形，称为\textbf{光滑带边区域}。其流形边界恰是它在 $M$ 中的拓扑边界。闭球是典型例子，而任意闭集未必是带边区域。
\end{example}

\begin{figure}[h!]
    \centering
\tikzset{
pattern size/.store in=\mcSize, 
pattern size = 5pt,
pattern thickness/.store in=\mcThickness, 
pattern thickness = 0.3pt,
pattern radius/.store in=\mcRadius, 
pattern radius = 1pt}
\makeatletter
\pgfutil@ifundefined{pgf@pattern@name@_hask52a3i}{
\pgfdeclarepatternformonly[\mcThickness,\mcSize]{_hask52a3i}
{\pgfqpoint{0pt}{0pt}}
{\pgfpoint{\mcSize+\mcThickness}{\mcSize+\mcThickness}}
{\pgfpoint{\mcSize}{\mcSize}}
{
\pgfsetcolor{\tikz@pattern@color}
\pgfsetlinewidth{\mcThickness}
\pgfpathmoveto{\pgfqpoint{0pt}{0pt}}
\pgfpathlineto{\pgfpoint{\mcSize+\mcThickness}{\mcSize+\mcThickness}}
\pgfusepath{stroke}
}}
\makeatother

% Pattern Info
 
\tikzset{
pattern size/.store in=\mcSize, 
pattern size = 5pt,
pattern thickness/.store in=\mcThickness, 
pattern thickness = 0.3pt,
pattern radius/.store in=\mcRadius, 
pattern radius = 1pt}
\makeatletter
\pgfutil@ifundefined{pgf@pattern@name@_9v4tzcn4d}{
\pgfdeclarepatternformonly[\mcThickness,\mcSize]{_9v4tzcn4d}
{\pgfqpoint{0pt}{0pt}}
{\pgfpoint{\mcSize+\mcThickness}{\mcSize+\mcThickness}}
{\pgfpoint{\mcSize}{\mcSize}}
{
\pgfsetcolor{\tikz@pattern@color}
\pgfsetlinewidth{\mcThickness}
\pgfpathmoveto{\pgfqpoint{0pt}{0pt}}
\pgfpathlineto{\pgfpoint{\mcSize+\mcThickness}{\mcSize+\mcThickness}}
\pgfusepath{stroke}
}}
\makeatother
\tikzset{every picture/.style={line width=0.75pt}} %set default line width to 0.75pt        

\begin{tikzpicture}[x=0.75pt,y=0.75pt,yscale=1,xscale=1]
%uncomment if require: \path (0,300); %set diagram left start at 0, and has height of 300

%Shape: Polygon Curved [id:ds6324304927731197] 
\draw   (14,85) .. controls (10,122) and (65.25,88.75) .. (74.25,80.25) .. controls (83.25,71.75) and (88.25,58.25) .. (81.25,45.25) .. controls (74.25,32.25) and (36.25,-5.25) .. (16.25,24.75) .. controls (-3.75,54.75) and (18,48) .. (14,85) -- cycle ;
%Shape: Rectangle [id:dp42790425123431874] 
\draw  [dash pattern={on 4.5pt off 4.5pt}] (142,117) node[below right]{$\mathbb H^n$} -- (252,117)node[below right]{$\R^n$} -- (252,7) -- (142,7) -- cycle ;
%Straight Lines [id:da2780621449052143] 
\draw    (197,117) -- (197,7) ;
%Curve Lines [id:da11211326883897499] 
\draw [pattern=_hask52a3i,pattern size=6pt,pattern thickness=0.75pt,pattern radius=0pt, pattern color={rgb, 255:red, 0; green, 0; blue, 0}] [dash pattern={on 4.5pt off 4.5pt}]  (197,25) .. controls (167,52) and (161,80) .. (196,97) ;
%Shape: Polygon Curved [id:ds6020575675206878] 
\draw  [pattern=_9v4tzcn4d,pattern size=6pt,pattern thickness=0.75pt,pattern radius=0pt, pattern color={rgb, 255:red, 0; green, 0; blue, 0}][dash pattern={on 4.5pt off 4.5pt}] (57.25,56.75) .. controls (55.75,74.75) and (65.25,88.75) .. (74.25,80.25) .. controls (83.25,71.75) and (88.25,58.25) .. (81.25,45.25) .. controls (74.25,32.25) and (58.75,38.75) .. (57.25,56.75) -- cycle ;
%Curve Lines [id:da8087172400108391] 
\draw [line width=0.75]    (77.25,67.75) .. controls (119.39,104.75) and (169.21,96.84) .. (192.37,74.39) node[near start,sloped,above]{$\psi$};
\draw [shift={(193.75,73)}, rotate = 133.75] [fill={rgb, 255:red, 0; green, 0; blue, 0 }  ][line width=0.08]  [draw opacity=0] (8.4,-2.1) -- (0,0) -- (8.4,2.1) -- cycle    ;
\end{tikzpicture}
    \caption{带边区域}
\end{figure}

最后确定边界的定向。外向是相对于 $M$ 的局部内侧而言，并不要求 $\partial M$ 将某个背景空间分成内外两部分。对 $p\in\partial M$，若 $\nu\in T_pM$ 满足 $\d x^n(\nu)<0$，就称 $\nu$ \textbf{外向}。此条件与边界坐标的选择无关：坐标变换保持边界和内侧，其微分在横向上只能乘以正数。这里不要求 $\nu$ 为单位矢量或与边界正交，因而无需度规。

\begin{theorem}\label{thm:boundary-orientation}
    定向带边流形 $M$ 自然诱导 $\partial M$ 的定向。若 $\epsilon$ 为 $M$ 的正定体元、$\nu$ 为沿边界的光滑外向矢量场、$\iota:\partial M\hookrightarrow M$ 为包含映射，则边界定向由
    \eq{\epsilon_{\partial}:=\iota^*(i_\nu\epsilon)}
    给出。对 $n\geqslant2$，这等价于规定：$(v_1,\cdots,v_{n-1})$ 是边界的右手基，当且仅当 $(\nu,v_1,\cdots,v_{n-1})$ 是 $M$ 的右手基，称为\textbf{外向在前}约定。
\end{theorem}
\begin{proof}
    各边界坐标系中的 $-\p_n$ 外向，用单位分解加权相加即可得到光滑外向场。因 $\nu$ 不在 $T_p(\partial M)$ 中，$\epsilon_{\partial}$ 处处非零。

    若另取外向场 $\nu'$，则沿边界有 $\nu'=a\nu+w$，其中 $a>0$、$w$ 与边界相切。由于 $n$ 个边界切矢线性相关，$\iota^*(i_w\epsilon)=0$，故 $\iota^*(i_{\nu'}\epsilon)=a\epsilon_{\partial}$。改变正定体元也只会乘以正函数，故所得定向无关这些选择。
\end{proof}

\begin{example}\label{eg:half-space-orientation}
    $\mathbb H^n$ 取标准体元 $\epsilon=\d x^1\wedge\cdots\wedge\d x^n$，外向矢量可取 $-\p_n$，于是
    \eq{\epsilon_{\partial}=(-1)^n\d x^1\wedge\cdots\wedge\d x^{n-1}.}
    因此 $n\geqslant2$ 时，边界坐标 $(x^1,\cdots,x^{n-1})$ 在 $n$ 为偶数时是右手系，为奇数时是左手系。
\end{example}

\begin{example}
    平面圆盘取 $\d x\wedge\d y$ 定向，其边界诱导逆时针方向。区间 $[a,b]$ 取 $\d x$ 定向，则左端点外向矢量为 $-\p_x$，右端点为 $\p_x$，诱导的零维体元分别为 $-1,+1$。零维流形的定向就是逐点指定正负号，故边界积分为 $\int_{\partial[a,b]}f=f(b)-f(a)$。
\end{example}

\section{广义 Stokes 定理}

考虑 $n$ 维定向流形 $M$。取连通坐标域上的坐标系 $(O,\psi)$，其坐标函数为 $x^i:O\to\R$。若右端 Riemann 积分存在，则定义 $A\subset O$ 上 $n$-形式 $\omega$ 的积分为
\eq{
    \int_A\omega:=\pm\int_{\psi[A]}(\omega_{1 \cdots n}\circ\psi^{-1})(x)\,\d x^1 \cdots\d x^n,
}
其中 $\pm$ 在右手系时取正。由多重积分的变量替换定理、$n$-形式分量变化律知，定义无关坐标选取，但保证 $A\subset O\cap O'$。
欲将积分域推至 $M$，需要将各局部积分进行缝合。已知单位分解 $\{\rho_i\}$ 满足 $\sum_i\rho_i=1$，因此 $\omega=\omega\sum_i\rho_i=\sum_i\rho_i\omega$。因单个 $\int\rho_i\omega$ 可用坐标系计算，故自然希望整个积分 $\int\omega$ 是 $\sum_i\int\rho_i\omega$。

\begin{definition}
    设 $\omega$ 为 $n$ 维定向流形 $M$ 上的紧支撑光滑 $n$-形式，取从属于坐标开覆盖 $\{O_i\}$ 的局部有限光滑单位分解 $\{\rho_i\}$，定义
    \eq{
        \int_M\omega:=\sum_i\int_{O_i}\rho_i\omega.
    }
    由紧支撑性及局部有限性，右端只有有限项非零。特别地，$M$ 紧时，任意光滑 $n$-形式均可如此积分。
\end{definition}
\begin{remark}
    实际上，总默认对诱导形式积分：$\int_S\omega=\int_S\tilde\omega$。
    本书只对阶数与积分域维数相同的形式定义积分；若将积分扩展到任意阶形式，可约定其余阶数的积分为零。
\end{remark}
\begin{theorem}
    上述积分与坐标开覆盖及从属单位分解的选取无关。
\end{theorem}
\begin{proof}
    设 $\{\rho_i\}$、$\{\rho'_j\}$ 分别从属于坐标开覆盖 $\{O_i\}$、$\{O'_j\}$。它们的支撑各只有有限个与紧集 $\supp\omega$ 相交，故以下均为有限和。又 $\rho_i\rho'_j\omega$ 的支撑含于 $O_i\cap O'_j$，且局部积分与坐标选择无关，于是
    \begin{align*}
        \sum_i\int_{O_i}\rho_i\omega
        &=\sum_{i,j}\int_{O_i\cap O'_j}\rho_i\rho'_j\omega\\
        &=\sum_j\int_{O'_j}\Bigl(\sum_i\rho_i\Bigr)\rho'_j\omega
        =\sum_j\int_{O'_j}\rho'_j\omega.\qedhere
    \end{align*}
\end{proof}

\begin{theorem}[广义 Stokes 定理]
    设 $\Omega$ 为 $n$ 维定向带边流形（$n\geqslant1$），$\partial\Omega$ 取诱导定向。对 $\Omega$ 上任意紧支撑光滑 $(n-1)$-形式 $\omega$，有
    \eq{\int_\Omega\d\omega=\int_{\partial\Omega}\omega.}
    特别地，$\Omega$ 紧时，任意光滑 $(n-1)$-形式均满足此式。
\end{theorem}
\begin{remark}
    该式统一了 Newton-Leibniz 定理、Green 公式、Gauss 定理和 Stokes 公式，将内部的微分与边界上的积分联系起来。
\end{remark}
\begin{proof}
    先在取标准定向的 $\mathbb H^n$ 上证明。写成
    \eq{\omega=\sum_{i=1}^n\omega_i\,\d x^1\wedge\cdots\wedge\widehat{\d x^i}\wedge\cdots\wedge\d x^n,}
    帽号表示略去该项，故
    \eq{\d\omega=\sum_{i=1}^n(-1)^{i-1}\pdv{\omega_i}{x^i}\,\d x^1\wedge\cdots\wedge\d x^n.}
    由紧支撑性，按 $x^i$ 积分时，$i<n$ 的项两端均为零；$i=n$ 的项只留下 $x^n=0$ 处的值。因此由 Fubini 定理及 Newton-Leibniz 公式，
    \eq{\int_{\mathbb H^n}\d\omega
        =(-1)^n\int_{\R^{n-1}}\omega_n(x^1,\cdots,x^{n-1},0)\,\d x^1\cdots\d x^{n-1}
        =\int_{\partial\mathbb H^n}\omega.}
    最后一步用了例 \ref{eg:half-space-orientation} 的边界定向，以及 $\d x^n$ 在边界上的拉回为零。对 $\R^n$，各项积分均为零，结论同样成立。

    一般地，取由连通坐标域组成的开覆盖 $\{O_i\}$ 及从属的局部有限光滑单位分解 $\{\rho_i\}$。各 $\rho_i\omega$ 紧支撑于 $O_i$，用坐标表示并在像域外延拓为零，即可应用上述结论；坐标系若反向，两边同时变号。又因只有有限项非零，且 $\sum_i\rho_i\omega=\omega$，故
    \[\int_\Omega\d\omega
        =\sum_i\int_\Omega\d(\rho_i\omega)
        =\sum_i\int_{\partial\Omega}\rho_i\omega
        =\int_{\partial\Omega}\omega.\qedhere\]
\end{proof}
\begin{remark}
    计算 $\int_\Omega\d\omega$ 时，可在各坐标域中略去边界，因为超平面 $x^n=0$ 对 $n$ 重 Riemann 积分无贡献。若将积分域记为 $\Omega^\circ$，应按上述局部积分理解；$\d\omega$ 限制到内部后未必仍有紧支撑。

    有的书会用 $\d^2=0$ 解释 $\partial(\partial\Omega)=\mt$，但这不能直接作为证明：前者是形式的微分恒等式，后者是边界的几何性质。若连续两次应用 Stokes 定理，须先确认 $\partial\Omega$ 的流形结构，还须说明对所有测试形式积分为零何以推出边界为空。本书已由半空间坐标直接证明后者，见定理 \ref{thm:boundary-no-boundary}。
\end{remark}

\begin{example}[Green]
    设 $A$ 是 $(\R^2,\delta)$ 上的对偶矢量场，$S$ 是单连通开集，则
    \eq{
    \int_S\left(\pdv{A_2}{x}-\pdv{A_1}{y}\right)\d x\,\d y=\oint_{\partial S} A_1\, \d x+A_2 \,\d y=\oint_{\partial S} A_\ell\,\d\ell.
    }
\end{example}
\begin{proof}
    置 1-形式 $\omega=A_\mu\mathrm{d} x^\mu$, 则
\[\mathrm{d} \omega  =\frac{\partial A_\mu}{\partial x^\nu} \mathrm{d} x^\nu \wedge \mathrm{d} x^\mu=\left(\frac{\partial A_2}{\partial x}-\frac{\partial A_1}{\partial y}\right) \mathrm{d} x \wedge \mathrm{d} y,\]
第一等式成立。
选线长 $\ell$ 为 $L$ 的局部坐标，把 $\tilde{\omega}$ 用坐标基展为 $\tilde{\omega}=\tilde{\omega}_1(\ell)\mathrm{d} \ell$, 两边作用于 $\partial / \partial \ell$ 得
$\tilde{\omega}_1(\ell)=\tilde{\omega}(\partial / \partial \ell)=\omega(\partial / \partial \ell)=A(\partial / \partial \ell)=A_\ell$，
故 $\tilde{\omega}=A_\ell\, \mathrm{d} \ell$。
\end{proof}
\begin{remark}
    $\int_{\partial S} \tilde{\omega}$ 的积分域是闭合曲线 $\partial S$，至少要用两个坐标域覆盖，因此应先对每个坐标域做局部积分再“缝合”。幸好可这样：设 $L$ 是 $\partial S$ 挖去一点所得，则它用一个坐标域即可覆盖，但挖去一点不影响 Riemann 积分。
\end{remark}

\section{Gauss 定理}\label{sec:gauss-divergence}

以下置 $\dim M=n\geqslant2$。沿用附录 \ref{appx:LA} 的定义，流形上的\textbf{度规}是逐点对称、非退化的光滑 $(0,2)$-张量场 $g$。本节允许不定度规，升降指标及正交性均相对于 $g$ 而言。

设 $M$ 已定向。由 \S\ref{sec:tensor-density}，度规与定向唯一确定\textbf{适配体元} $\epsilon$；在右手坐标系中，
\eq{\label{eq:metric-volume}
    \epsilon=\sqrt{|g|}\,\d x^1\wedge\cdots\wedge\d x^n,
    \qquad g:=\det(g_{\mu\nu}).
}
对紧支撑连续函数 $f$，沿用前节的局部积分及单位分解，约定 $\int_M f:=\int_M f\epsilon$。例如，时空作用量 $S=\int_M\mathcal L$ 省略了适配体元 $\sqrt{-g}\,\d[4]x$。

\begin{definition}
    $M$ 的余维为 $1$ 的嵌入子流形 $\Sigma$ 称为\textbf{超曲面}(hypersurface)。记包含映射为 $\iota:\Sigma\hookrightarrow M$，通过 $\iota_*$ 将 $T_p\Sigma$ 视为 $T_pM$ 的子空间。满足 $g(\nu,w)=0$（$\forall w\in T_p\Sigma$）的非零切矢 $\nu\in T_pM$ 称为 $\Sigma$ 在 $p$ 处的\textbf{法矢}。
\end{definition}

超曲面局部可写为 $\Sigma=\{f=0\}$，其中 $\d f\ne0$。此时
\eq{T_p\Sigma=\ker(\d f)_p,\qquad
    (T_p\Sigma)^\perp=\operatorname{span}\{(\d f)_p^\sharp\}.}
因为消去 $T_p\Sigma$ 的余矢构成一维空间，而音乐同构将其映到法空间，故任意两法矢只差非零倍数。法矢场局部总存在，整体存在则还需能一致地选取法向。

\begin{definition}
    $\Sigma$ 上的\textbf{诱导度规}为 $\tilde g:=\iota^*g$，即只让 $g$ 作用于超曲面的切矢。若以 $y^a$ 为 $\Sigma$ 的坐标、嵌入写作 $x^\mu=x^\mu(y)$，则
    \eq{\tilde g_{ab}=g_{\mu\nu}\pdv{x^\mu}{y^a}\pdv{x^\nu}{y^b}.}
\end{definition}

诱导度规未必非退化：其退化子空间为 $T_p\Sigma\cap(T_p\Sigma)^\perp$，而法空间只有一维，故它退化当且仅当法矢满足 $g(\nu,\nu)=0$。此时超曲面在该点\textbf{类光}，法矢本身也切于 $\Sigma$。以下法向分解均假设 $\Sigma$ 非类光，并取单位法矢
\eq{\sigma:=g(\nu,\nu)=\pm1.}
每个 $v\in T_pM$ 都有唯一分解
\eq{v=\sigma g(v,\nu)\nu+h(v),\qquad h(v)\in T_p\Sigma.}
于是切向\textbf{投影映射}及其降指标形式为
\eq{h^\mu{}_{\lambda}=\delta^\mu{}_{\lambda}-\sigma\nu^\mu\nu_{\lambda},
    \qquad h_{\mu\lambda}=g_{\mu\lambda}-\sigma\nu_\mu\nu_\lambda.}
显然 $h^2=h$、$h(\nu)=0$，且 $h$ 在 $T_p\Sigma$ 上为恒等映射。$h_{\mu\lambda}$ 在整个 $T_pM$ 上退化，限制到 $T_p\Sigma$ 才是非退化的 $\tilde g$；它是满足 $h(\nu,\cdot)=0$ 的诱导度规的唯一对称延拓。

\begin{remark}
    适配坐标 $\Sigma=\{x^n=0\}$ 中，$\d x^n$ 是法余矢，法矢则正比于 $(\d x^n)^\sharp$，一般并非 $\p_n$。
\end{remark}

\begin{theorem}
    设 $\Sigma$ 非类光，并选定光滑单位法矢场 $\nu$。以 $(\nu,v_1,\cdots,v_{n-1})$ 为右手基来规定 $\Sigma$ 的定向，则其适配体元为
    \eq{\tilde\epsilon=\iota^*(i_\nu\epsilon).}
    特别地，$\Sigma=\partial D$ 且 $\nu$ 外向时，此定向就是前文的边界诱导定向。
\end{theorem}
\begin{proof}
    逐点取正交归一右手基 $(\nu,e_1,\cdots,e_{n-1})$，其中 $e_a$ 切于 $\Sigma$。则 $e_a$ 关于 $\tilde g$ 也正交归一，且
    \eq{\tilde\epsilon(e_1,\cdots,e_{n-1})
        =\epsilon(\nu,e_1,\cdots,e_{n-1})=1,}
    故 $\tilde\epsilon$ 正是该定向下的适配体元。
\end{proof}

在 $\Sigma$ 的右手坐标系中，$\tilde\epsilon=\sqrt{|\det(\tilde g_{ab})|}\,\d y^1\wedge\cdots\wedge\d y^{n-1}$。注意，$\iota^*\epsilon=0$，因为 $\Sigma$ 只有 $n-1$ 维；诱导体元须先对法矢内乘，再拉回。

\begin{example}
    欧氏空间的球面 $r=R$ 取外法矢 $\nu=\p_r$。由
    \eq{g=\d r^2+r^2\d\theta^2+r^2\sin^2\theta\,\d\phi^2,\qquad
        \epsilon=r^2\sin\theta\,\d r\wedge\d\theta\wedge\d\phi,}
    得 $\tilde g=R^2(\d\theta^2+\sin^2\theta\,\d\phi^2)$、$\tilde\epsilon=R^2\sin\theta\,\d\theta\wedge\d\phi$，即熟悉的球面度规与面积元。
\end{example}

\begin{definition}
    矢量场 $X$ 相对于 $g$ 的\textbf{散度} $\operatorname{div}_gX$ 由
    \eq{\Lie_X\epsilon=(\operatorname{div}_gX)\epsilon}
    定义，即体元沿 $X$ 的相对变化率。
\end{definition}

\begin{theorem}
    散度满足
    \eq{\label{eq:divergence-volume}
        \d(i_X\epsilon)=(\operatorname{div}_gX)\epsilon,\qquad
        \operatorname{div}_gX=\frac{1}{\sqrt{|g|}}\p_\mu\bigl(\sqrt{|g|}X^\mu\bigr).
    }
\end{theorem}
\begin{proof}
    因 $\epsilon$ 为最高阶形式，$\d\epsilon=0$，第一式即 Cartan 公式。对式 \eqref{eq:metric-volume} 内乘后取外微分，得
    \eq{\d(i_X\epsilon)
        =\p_\mu\bigl(\sqrt{|g|}X^\mu\bigr)\,\d x^1\wedge\cdots\wedge\d x^n,}
    与第一式比较即得第二式。
\end{proof}

\begin{remark}
    引入度规的 Levi-Civita 联络后（见附录 \ref{appx:curvature}），此散度就是 $\nabla_\mu X^\mu$。因为
    \eq{\Gamma^\mu{}_{\mu\lambda}
        =\frac12g^{\alpha\beta}\p_\lambda g_{\alpha\beta}
        =\p_\lambda\ln\sqrt{|g|},}
    故 $\nabla_\mu X^\mu=\p_\mu X^\mu+\Gamma^\mu{}_{\mu\lambda}X^\lambda$ 恰为上式。这里不能将 Levi-Civita 联络换为任意仿射联络。
\end{remark}

\begin{theorem}[Gauss]\label{thm:gauss-divergence}
    设 $D$ 为定向度规流形 $(M,g)$ 中的光滑带边区域，$X$ 为光滑矢量场，且其在 $D$ 上的限制紧支撑。记 $j:\partial D\hookrightarrow M$，则
    \eq{\int_D(\operatorname{div}_gX)\epsilon
        =\int_{\partial D}j^*(i_X\epsilon).}
    若边界处处非类光，取单位外法矢 $\nu$、$\sigma=g(\nu,\nu)$ 及 $\tilde\epsilon=j^*(i_\nu\epsilon)$，则又有
    \eq{\int_D(\operatorname{div}_gX)\epsilon
        =\int_{\partial D}\sigma g(X,\nu)\tilde\epsilon.}
    $D$ 紧时，无需另加紧支撑条件。
\end{theorem}
\begin{proof}
    将式 \eqref{eq:divergence-volume} 代入 Stokes 定理即得第一式。对第二式，在边界分解 $X=\sigma g(X,\nu)\nu+h(X)$。因 $h(X)$ 切于边界，$\epsilon$ 作用于它及其余 $n-1$ 个边界切矢时为零，故
    \eq{j^*(i_X\epsilon)=\sigma g(X,\nu)\tilde\epsilon.\qedhere}
\end{proof}

\begin{remark}
    正定度规下 $\sigma=1$，即通常的散度定理；按本书的 Lorentz 号差约定，类时单位法矢满足 $\sigma=-1$，不可略去此号。类光边界没有单位法矢，不能使用第二式，但第一式仍成立。
\end{remark}

为与 \S\ref{sec:Gauss} 的分量写法对应，定义\textbf{有向面元}的分量
\eq{\label{eq:directed-surface-element}
    \d\Sigma_\mu:=j^*(i_{\p_\mu}\epsilon)
        =\frac{1}{(n-1)!}\epsilon_{\mu\mu_2\cdots\mu_n}
        j^*(\d x^{\mu_2}\wedge\cdots\wedge\d x^{\mu_n}).
}
于是 $j^*(i_X\epsilon)=X^\mu\d\Sigma_\mu$；非类光边界上还有 $\d\Sigma_\mu=\sigma\nu_\mu\tilde\epsilon$。Gauss 定理在坐标域内写成
\eq{\int_D\p_\mu\bigl(\sqrt{|g|}X^\mu\bigr)\,\d[n]x
    =\int_{\partial D}X^\mu\d\Sigma_\mu,}
其中体积分采用右手坐标，一般情形按坐标分片求和。四维 Lorentz 度规下 $\sqrt{|g|}=\sqrt{-g}$，这就是式 \eqref{eq:Gauss}。

\section{Frobenius 定理}\label{sec:Frobenius}

指定 n 维流形的 m 维子空间
⼦空间场（m 维分布）、光滑 m 维分布定义

积分子流形、可积的、最大的

可积充要条件：
Frobenius 定理（矢量表述）

%当条件满足时，不但m维分布可积，而且 \forall p\in M 都存在唯一的最大积分子流形\Sigma 使%p\in\Sigma。这是 Frobenius 定理的整体表述中的一个结论。证明:\dots

%example: 下面给出一个用 Frobenius 定理求解偏微分方程组的例子。
%假定要求解关于未知函数u(x,y)的下列偏微分方程组

%设W是n维流形M上的一个n-1维分布.满足如下条件的余矢场称为消灭（annihilated）余矢场
%...
%可见任一p\in M的 T_p^*M 中的所有消灭余矢构成 T_p^*M 的一个1维子空间

%theorem
%n维流形M上的n-1维分布W 可积的充要条件是：对W的任一可微的消灭余矢场 \omega 及任意 p\in M，存在 p 的%开邻域 U及 W在 U上的消灭余矢场 \mu 和 上的余矢场 \psi 使
%\d\omega = \mu \wedge \psi
%proof

\begin{example}
    广义相对论中，超曲面正交(hypersurface orthogonal)，因此无旋参考系又称为超曲面正交参考系。有旋就相应于非超曲面正交。
\end{example}
\begin{proof}
    首先我们要先明确“超曲面正交”是何含义：对参考系定义域内的任一点 $p$，存在过$p$的一张超曲面 $\Sigma$ 同系内所有观者世界线正交。以 $Z$ 代表参考系的 4-速场，$Z_\flat=g_\flat(Z)$，则接受微分形式的结论：以上表述等价于
    \eq{
    Z\wedge\d Z=\A(Z\otimes\nabla Z)=0,
    }
    这里 $\d Z$ 表示 $Z$ 的外微分。
\end{proof}

\section{积分的求导}

\begin{theorem}
    设完备向量场 $v\in\mathfrak X(\R^n)$作用下 $p$维带边闭子流形 $D_0\subseteq\R^n$随时间演化为 $D\subseteq\R^{n+1}$，对同样随时演化的微分形式 $\omega\in\Omega^p(\R^n)\times\R$有 
 \eq{
    \dv{t}\int_{D_t}\omega=\int_{D_t}\frac{\partial\omega}{\partial t}+\int_{D_t}i_v\d\omega+\int_{\partial D_t}i_v\omega,
 }
\end{theorem}
\begin{proof}
    设 $v$由单参数群 $\psi$给出，设 $\R^{n+1}$上的向量场 $X:=v+\partial_t$由单参数群 $\phi$给出。注意到 $\omega$不含 $\d t$有
    \eq{
        i_X\omega=i_v\omega,\quad i_X\d\omega=i_v\d_x\omega+i_{\partial_t}(\dot{\omega}_I\d t\wedge\d x^I)=i_v\d_x\omega+\dot\omega.
    }
    则依 Lie导数定义、Cartan公式、Stokes公式有
    \begin{align*} \dv{t}\bigg|_{t=0}\int_{D_t}\omega&=\dv{t}\bigg|_{t=0}\int_{D_0}\phi_t^*\omega =\int_{D_0}\dv{t}\bigg|_{t=0}\phi_t^*\omega=\int_{D_0}\mathcal L_X\omega\\ &=\int_{D_0}\d i_X\omega+i_X\d\omega =\int_{\partial D_0}i_v\omega+\int_{D_0}i_v\d_x\omega+\int_{D_0}\frac{\partial\omega}{\partial t}. \qedhere \end{align*}
\end{proof}

\begin{example}
在三维注意到 $2$-形式 $\omega$ 和向量场 $v$ 满足
\begin{align*}
    \star(\star\omega\wedge v^\flat)&=\star(\epsilon^{ijk}\omega_{ij}g_{k\mu}v_\nu\mathrm dx^\mu\wedge\mathrm dx^\nu)\\
    &=\epsilon_{\alpha\beta\gamma}\epsilon^{ijk}\omega_{ij}v_\nu g_{k\mu}g^{\mu\alpha}g^{\nu\beta}\mathrm dx^\gamma\\ 
    &=\delta_{\beta\gamma k}^{ijk}\omega_{ij}v^\beta\mathrm dx^\gamma\\ 
    &=(\omega_{ij}-\omega_{ji})v^i\mathrm dx^j=\iota_v\omega, 
\end{align*}

在 $\R^3$ 矢量代数中，我们有
令 $\omega=\mathbf F^\flat$ 就有 $\iota_v\omega=\mathbf F\cdot\bm v,\quad \iota_v\mathrm d\omega=\star(\star\mathrm d\mathbf{F}^\flat\wedge\bm v^\flat)=((\nabla\times\mathbf F)\times \bm v)^\flat. $
令 $\omega=\star\mathbf{F}^\flat$ 就有 $\iota_v\omega=\star(\mathbf{F}^\flat\wedge\bm{v}^\flat)=(\mathbf F\times\bm v)^\flat,\quad\iota_v\mathrm d\omega=\star((\star\mathrm d\star\mathbf{F}^\flat)\bm v^\flat)=\star((\nabla\cdot\mathbf F)\bm v)^\flat. $
令 $\omega=\star\mathrm F$ 就有 $\iota_v\omega=\star(\mathrm F\bm v)^\flat,\quad\iota_v\mathrm d\omega=\star(\star\mathrm d\star\mathrm F\wedge\bm v^\flat)=0.  $
\end{example}
